\section{Lifetime timeline}\label{sec:lifetime}
The timeline runs from formation to the conditional disappearance of the
remnant. Computed entries are labelled by track and clock; the rest are
literature values or conditional stages whose ages remain to be calculated.
Several late processes can overlap, and disappearance means loss of the
baryonic remnant through decay products.

{\small\renewcommand{\arraystretch}{1.12}
\begin{longtable}{@{}>{\raggedright\arraybackslash}p{0.17\textwidth}>{\raggedright\arraybackslash}p{0.22\textwidth}>{\raggedright\arraybackslash}p{0.56\textwidth}@{}}
\toprule Age or condition & Event & Basis \\\midrule\endfirsthead
\toprule Age or condition & Event & Basis \\\midrule\endhead
Before $t=0$ & Formation and accretion & Cloud collapse, disk and outflows set the mass, entropy and angular momentum; outside the fixed-mass calculation. Low-mass pre-main-sequence models provide the comparison \cite{baraffe15}. \\[3pt]
Contraction clock, 0--1 Gyr & Deuterium burning and contraction & Computed: \pmsTemperature\ K and $\pmsRadius\Rs$ near 1 Gyr with \pmsNuclearPercent\% nuclear support; continued to $\continuousAge\Tyr$ (Section~\ref{sec:contraction}). Lithium depletion is not followed. \\[3pt]
Long-term clock, $t=0$ & Static hydrogen-burning model & Tracks S, F and P start here with $X=0.7$, $\Teff=2766$ K. \\[3pt]
$0.7103\Tyr$ (P) & Maximum $^3$He & Central mass fraction 0.1003, then destruction overtakes production. LBA97: 0.0995 at $1.38\Tyr$. \\[3pt]
$3.546$--$3.551\Tyr$ (P, F, S) & Loss of full convection & A radiative layer separates centre and envelope at central $X\simeq0.175$; the centre becomes radiative near $3.551\Tyr$ in F and S. LBA97: $5.742\Tyr$, $X\simeq0.16$. \\[3pt]
$3.650$--$3.767\Tyr$ (S) & Central hydrogen depletion & Central $X$ crosses $10^{-3}$, $10^{-4}$ and $10^{-5}$; $0.0147\Ms$ of hydrogen remains at the first crossing and burns outside the centre. \\[3pt]
$3.673$, $3.838\Tyr$ (S) & Luminosity and temperature maxima & $0.006513\Ls$; 4754 K. Burning still dominates; no cooling sequence is established. \\[3pt]
$4.002\Tyr$ (S) & Last computed model & 4612 K, central $X=3.528\times10^{-8}$, 97.28\% nuclear. \\[3pt]
Provisional & Helium-3 shell flash & Occurs after a boundary switch; dependence on the original atmosphere untested (Section~\ref{sec:flash}). \\[3pt]
To be calculated & Residual burning fades & Record nuclear heat below 10\%, 1\% and 0.1\% of the photon luminosity, including any recovery. LBA97 reaches 5807 K at its turn; a comparison value only. \\[3pt]
To be calculated & Diffusion and separation & Hydrogen collects at the surface while heavier species settle; a helium core does not imply a helium atmosphere. \\[3pt]
To be calculated & Molecules, grains, dense atmospheres & Follow molecular and collision-induced absorption, condensation and settling where local conditions permit \cite{chabrier00,blouin18}. \\[3pt]
To be calculated & Crystallization and separation & Helium and mixture phase equilibria, latent heat, solid heat capacity; the front follows density and core temperature, not surface temperature \cite{baiko22}. \\[3pt]
To be calculated & Envelope and core couple & Whether and when the convective envelope reaches the nearly isothermal interior; order relative to crystallization unknown. \\[3pt]
$\Teff=1000$, 500, 100 K & Cold-remnant milestones & Record age, core temperature, radius, composition, solid fraction and power sources; below 100 K continue with a solid, partly neutral description. \\[3pt]
$\gtrsim10^{14}$ yr & Changing environment & Retained and ejected objects in evolving stellar and dark-matter environments; accretion of gas, debris and solid bodies is conditional; no unique encounter rate is assumed. \\[3pt]
Close passages & Tidal heating, capture, collision & Distinguish heat from spin and oscillations; couple an encounter survey to waiting-time distributions and post-pulse cooling. \\[3pt]
Allowed interactions & Dark-matter heating & Compare capture, thermalization, evaporation and annihilation with the photon luminosity; include a no-annihilation case \cite{bell24}. \\[3pt]
Nuclear times competitive & Density-driven fusion & Pycnuclear and impurity reactions against cooling and decay times in the helium-rich composition \cite{gasques05}. \\
\bottomrule
\end{longtable}}

\paragraph{Decay and disappearance (conditional).}
Proton decay has not been detected; the Super-Kamiokande partial-lifetime
limit for $p\rightarrow e^+\pi^0$ is $\tau_p/B>2.4\times10^{34}$ yr
\cite{takenaka20}, and baryogenesis does not require it \cite{fukugita86}.
Virtual black holes supply a possible channel with
$\tau_p\sim10^{45}$ yr \cite{adams01}, distinct from the far slower
collective tunnelling of degenerate matter \cite{adams98}. We therefore
carry finite-lifetime scenarios and a stable-baryon branch rather than one
disappearance date. In order of condition, the stages are: decay heating
exceeds cooling and environmental power, possibly long before much mass is
lost; accumulated decays change the composition and mass; degeneracy is
lost and the mass--radius relation must be recalculated for the actual
composition; decay products and then thermal radiation escape as columns
shrink, so that an optically thin object need not radiate as a blackbody;
the remnant ceases to be bound, which is distinct from the decay of every
former baryon; and finally discrete survival replaces a continuous mass. In
the stable-baryon branch cooling, nuclear and environmental alternatives
continue without a decay-imposed endpoint.

\paragraph{A reference clock for the decay era.}
If every baryon has the same constant, independent mean lifetime $\tau_B$,
with no accretion or stable daughters, and $t_d$ measures time from the
start of the comparison,
\begin{equation}
 \langle M_B(t_d)\rangle=M_{B,0}e^{-t_d/\tau_B},\qquad
 L_{\rm heat}=f_{\rm dep}(t_d)\frac{\langle M_B\rangle c^2}{\tau_B},
\end{equation}
where the deposited fraction $f_{\rm dep}$ may evolve and a physical channel
would also account for daughter masses and binding. Half the mass is lost at
$t_d=0.6931\,\tau_B$, 90\% at $2.303$, 99\% at $4.605$ and 99.99\% at
$9.210\,\tau_B$. For $N_0$ independently decaying baryons,
\begin{equation}
 \Pr(t_{\rm last}\leq t_d)=\left(1-e^{-t_d/\tau_B}\right)^{N_0},\qquad
 \langle t_{\rm last}\rangle=\tau_B H_{N_0},
\end{equation}
with $H_N$ the harmonic number. For $0.1\Ms$, $N_0\simeq1.2\times10^{56}$,
the mean last decay occurs near $129.7\,\tau_B$ with a central 90\% interval
of 128.0--132.1$\,\tau_B$. This dates the last of the original baryons,
dispersed or not, not the survival of a bound star. These are analytic
consequences of the stated assumptions, not an extrapolation of any
computed track; a reference lifetime is not a measured age, and changing it
changes temperature, chemistry and energy escape as well as the time scale.
